\documentclass[12pt, oneside]{book}
% This provides the \BibTeX macro
\usepackage{doc}
\usepackage{makeidx}

\usepackage[tight,footnotesize]{subfigure}
\usepackage{amsmath}
\usepackage{amsfonts}

\newtheorem{problem}{Problem}
\newtheorem{definition}{Definition}
\usepackage{graphicx}
\usepackage{algorithmic}

\begin{document}
\begin{center}
\begin{minipage}{0.75\linewidth}
    \centering
    \includegraphics[width=0.3\linewidth]{logo}
    \par
    \vspace{3cm}
    {\uppercase{\Large Preconditioning\par}}
    \vspace{3cm}
    {\Large Mohammad Ali Rostami\par}
    \vspace{3cm}
    {\Large Supervisor: Prof. Martin B{\"u}cker\par}
    \vspace{3cm}
    {\Large January 2016}
\end{minipage}
\end{center}
\clearpage

\newpage
\section*{Abstract}

\thispagestyle{empty}

\newpage
\section*{Acknowledgments}
\thispagestyle{empty}


\tableofcontents

\chapter{Introduction}
\chapter{Preliminaries}
\begin{definition}[Structurally Orthogonal]
Two columns (similary two rows) $c_i$ and $c_j$ are structurally orthogonal if and only if
they do not have any nonzero element in the same row (similarly in the same column).
\end{definition}

\begin{definition}[Partially Structurally Orthogonal]
Given the set of required elements $R$, two columns $c_i$ and $c_j$ are structurally orthogonal 
if and only if they do not have any nonzero element from $R$ in the same row.
\end{definition}



\section{Graph Models}
The graph model for ILU preconditioning can be expressed as the directed graph $G_F=(V,E)$
in which $V$ represents rows (and columns). The edges of this graph are actually the adjacency relation
and actually the nonzeros. (Fill path)

The graph model for partial one-soded coloring would be an undirected graph in which
vertices are defined as columns. Two vertices are connected if their corresponding 
columns have a nonzero elements in a same row. Being partial is considered due given
required elements as a pattern. 

Given the matrix $A$, the existing algorithm from [michael's thesis] can be expresses as follows,
\begin{itemize}
\item $R_{Init}=\rho(A)$
\item Preconditioning: $[A^{'},F] = SILU(R_{init},el)$
\item Coloring: $\Phi(R_{init},A)$
\item $R_{pot}\subset A\ R_{init}$ such that $|\Phi(R_{init},A)|=|\Phi(R_{init}\cup R_{pot},A)|$
\item $R_{add}\subset A\ R_{pot}$ such that $|SILU(R_{init}\cup R_{add},el)|=|SILU(R_{init},el) \cup R_{add}|$
\end{itemize}

We can introduce the new algorithm as follows.
\begin{itemize}
\item $R_{Init}=\rho(A)$
\item Find the graph models for coloring $G_{\Phi}$ and for preconditioning $G_{ILU}$
\item 
\begin{align}
&G_{ILU}^{0} \rightarrow G_{ILU}^{1}\\
&G_{\Phi}^{0} \rightarrow G_{\Phi}^{1}
\end{align}
\end{itemize}
The goal is to minimize the number of fillins as well as the number of colors in a way that
the number of additionally required elements is maximized.
%\chapter{Conjecture Checking on list of Graph}
%\chapter{Tools for Combinatorial Scientific Computing}
%\section Conjecture Checking with GraphTea
%\section EXPLAIN
\chapter{EXPLAIN}
\section{Column Compression}
\cite{2013:05,2014:01}
\section{Bidirectional Compression}
\cite{2014:09}
\section{Partial Jacobian Computation}
...
\section{Nested Dissection Ordering}
\cite{2014:02}
\section{Parallel Matrix Vector Product}
\cite{2015:3}
\section{EXPLAIN Revisited}
Changing the whole structure of back bone to speed up the software.
\chapter{GraphTea}
\cite{2014:07}
\cite{2014:15}
\cite{2014:16}

Chemical Graph theory
\cite{2015:05,2015:06,2015:07,2015:08}

\chapter{Conclusion}
 
\bibliographystyle{IEEEtran}
\bibliography{refs}

\end{document}
